\documentstyle[%


righttag%


,11pt%


,amssymb%


,russian%


,izvru%


]{amsart}





\newcommand{\tl}[3]{\medskip\noindent{\bf#1}\ #2 \dotfill #3 \par} 


\pagestyle{empty}


\begin{document}


%\tableofcontents


\par


\vskip 2cm


Данный тематический номер журнала посвящен математическим моделям


и численным методам. Составитель номера --- профессор 


Ляшко А.Д.


\par


\bigskip


\bigskip


\bigskip


\par


\centerline{\Large СОДЕРЖАНИЕ}


\bigskip


\bigskip





\tl{Амосов А.А., Злотник А.А.}


{Полудискретный метод решения уравнений одномерного \newline


движения неоднородного вязкого


теплопроводного газа с негладкими данными}{3}





\tl{Глазырина Л.Л., Павлова М.Ф.}


{О разрешимости одного нелинейного эволюционного\newline


неравенства теории совместного движения


поверхностных и подземных вод}{20}





\tl{Дьяконов Е.Г.}


{Оценки $N$--поперечников в смысле Колмогорова


для некоторых компактов\newline


в усиленных пространствах Соболева}{32}





\tl{Зайцева С.Б., Злотник А.А.}


{Точные оценки градиента погрешности


локально-\newline


одномерных методов для многомерного


уравнения теплопроводности}{51}





\tl{Карчевский М.М., Ляшко А.Д., Паймушин В.Н.}


{О математических задачах теории\newline


многослойных 


оболочек с трансверсально-мягкими заполнителями}{66}





\tl{Крукиер Л.А.}


{Кососимметричные итерационные методы решения 


стационарной задачи\newline


конвекции-диффузии 


с малым параметром при старшей производной}{77}





\tl{Федотов Е.М.}


{Об одном классе двухслойных разностных схем 


для нелинейных краевых\newline


задач с памятью}{86}





\tl{Шишкин Г.И., Целищева И.В.}


{Метод декомпозиции для сингулярно возмущенных\newline


краевых задач с локальным возмущением начальных 


условий. Уравнения с\newline


конвективными членами}{98}





\tl{Мартюшов С.Н.}


{Построение дву- и трехмерных сеток для задач 


газодинамики на основе\newline


уравнения Пуассона}{108}





\vfill


\noindent\raisebox{2pt}{\copyright} Казанский госудаpственный унивеpситет


\end{document}


